\suggestions

A partir de los resultados obtenidos y de las proyecciones tecnológicas identificadas durante el ciclo de desarrollo, se proponen las siguientes líneas de trabajo futuro:

Se sugiere expandir el banco de preguntas relacionadas con los Objetivos de Desarrollo Sostenible, mediante la incorporación de múltiples niveles de complejidad cognitiva que doten al sistema de una mayor adaptabilidad para interactuar con diferentes grupos etarios.

Resulta pertinente implementar un sistema de logros y recompensas virtuales que fortalezca el paradigma de gamificación, de modo que se logre un mayor incentivo para que los jugadores completen desafíos específicos y asimilen el contenido de sostenibilidad a largo plazo.

Se exhorta a diseñar una arquitectura de red para habilitar un modo multijugador en línea, lo cual permitiría la interconexión de estudiantes de diferentes instituciones escolares y promovería el intercambio dinámico de conocimientos.

Se recomienda enriquecer el diseño visual del proyecto mediante la adición de microanimaciones y efectos de transición de turnos elaborados, maximizando la inmersión de la experiencia de juego.

Finalmente, se propone refactorizar la vista y los controles del sistema para su distribución nativa en plataformas móviles, lo que permitiría ampliar drásticamente el alcance y la accesibilidad del recurso educativo en la población.
